\section{Results}\label{sec:results}
% Written in P9 (docs/DECISIONS.md K-069). Sources: results/tables/{stats,all_results,delta_f1,pareto,
% calibration_shift}.csv through generated/numbers.tex (analysis/paper_numbers.py: rq1_macros, rq2_macros,
% rq3_macros, calib_macros; the qualitative claims below are checked there and stop the generator if they fail).
% Decisions: K-036, K-041, K-043, K-044, K-050, K-051, K-057 ... K-066. No numeric literal: every number is a macro.
% Figure placement (docs/FIGURES.md, K-066): Fig. 3, 5, 8 here; Fig. 4, 7 and the F1 table in the appendix.

All results are on the test split. Unless stated otherwise, each quantized version uses its own threshold, fitted on
its validation scores at a false-positive rate of \TargetFPR{} (Section~\ref{sec:method}); the loss is
$\Delta = \Qfp{} - \mathrm{Q}k$, so a positive value is a loss. Values are means over the \NSeeds{} seeds, and brackets
are \CIPct{} hierarchical bootstrap intervals over seeds and test windows (Section~\ref{sec:setup}). Each hypothesis is
restated where it is tested; its verdict rule was fixed after the descriptive results of the full grid and before
any interval or test was computed.

\begin{table*}[t]
\centering
\footnotesize
\setlength{\tabcolsep}{3pt}
\caption{Test Recall (mean $\pm$ s.d. over 3 seeds) per attack type, model and quantization level (Q0 fp32, Q1 fp16, Q2 INT8); each version uses its own validation threshold (FPR = 1\%). XGBoost is deterministic over seeds (s.d. 0). $^\dagger$ test-only pair, outside the primary pool.}
\label{tab:recall}
\begin{tabular}{@{}llrrrrrrr@{}}
\toprule
& & \multicolumn{3}{c}{1D-CNN} & \multicolumn{3}{c}{conv-AE} & XGBoost \\
\cmidrule(lr){3-5}\cmidrule(lr){6-8}
Attack type & Group & Q0 & Q1 & Q2 & Q0 & Q1 & Q2 & (not quantized) \\
\midrule
\multicolumn{9}{l}{\textit{Car-Hacking}} \\
DoS & coarse & 1.000 $\pm$ 0.000 & 1.000 $\pm$ 0.000 & 1.000 $\pm$ 0.000 & 0.983 $\pm$ 0.003 & 0.983 $\pm$ 0.003 & 0.983 $\pm$ 0.004 & 1.000 $\pm$ 0.000 \\
fuzzy & coarse & 0.999 $\pm$ 0.000 & 0.999 $\pm$ 0.000 & 0.999 $\pm$ 0.000 & 0.996 $\pm$ 0.000 & 0.996 $\pm$ 0.000 & 0.996 $\pm$ 0.000 & 1.000 $\pm$ 0.000 \\
gear spoofing & coarse & 1.000 $\pm$ 0.000 & 1.000 $\pm$ 0.000 & 1.000 $\pm$ 0.000 & 0.874 $\pm$ 0.092 & 0.874 $\pm$ 0.093 & 0.899 $\pm$ 0.096 & 1.000 $\pm$ 0.000 \\
RPM spoofing & coarse & 1.000 $\pm$ 0.000 & 1.000 $\pm$ 0.000 & 1.000 $\pm$ 0.000 & 0.994 $\pm$ 0.001 & 0.994 $\pm$ 0.001 & 0.994 $\pm$ 0.001 & 1.000 $\pm$ 0.000 \\
pooled (primary pool) & -- & 1.000 $\pm$ 0.000 & 1.000 $\pm$ 0.000 & 1.000 $\pm$ 0.000 & 0.955 $\pm$ 0.030 & 0.955 $\pm$ 0.030 & 0.963 $\pm$ 0.031 & 1.000 $\pm$ 0.000 \\
\midrule
\multicolumn{9}{l}{\textit{ROAD}} \\
fuzzing & coarse & 0.995 $\pm$ 0.008 & 0.995 $\pm$ 0.008 & 1.000 $\pm$ 0.000 & 0.564 $\pm$ 0.221 & 0.569 $\pm$ 0.221 & 0.539 $\pm$ 0.234 & 0.779 $\pm$ 0.000 \\
correlated signal & coarse & 0.901 $\pm$ 0.082 & 0.901 $\pm$ 0.082 & 0.898 $\pm$ 0.085 & 0.036 $\pm$ 0.038 & 0.035 $\pm$ 0.036 & 0.044 $\pm$ 0.055 & 0.254 $\pm$ 0.000 \\
max speedometer & coarse & 0.642 $\pm$ 0.093 & 0.641 $\pm$ 0.093 & 0.635 $\pm$ 0.091 & 0.069 $\pm$ 0.061 & 0.069 $\pm$ 0.060 & 0.072 $\pm$ 0.067 & 0.486 $\pm$ 0.000 \\
reverse light off & coarse & 0.356 $\pm$ 0.115 & 0.356 $\pm$ 0.115 & 0.350 $\pm$ 0.110 & 0.000 $\pm$ 0.000 & 0.000 $\pm$ 0.000 & 0.000 $\pm$ 0.000 & 0.624 $\pm$ 0.000 \\
reverse light on & coarse & 0.088 $\pm$ 0.020 & 0.087 $\pm$ 0.020 & 0.083 $\pm$ 0.019 & 0.001 $\pm$ 0.001 & 0.001 $\pm$ 0.001 & 0.001 $\pm$ 0.001 & 0.215 $\pm$ 0.000 \\
correlated signal (masq.) & stealthy & 0.854 $\pm$ 0.064 & 0.854 $\pm$ 0.064 & 0.852 $\pm$ 0.061 & 0.021 $\pm$ 0.017 & 0.021 $\pm$ 0.017 & 0.025 $\pm$ 0.028 & 0.070 $\pm$ 0.000 \\
max speedometer (masq.) & stealthy & 0.718 $\pm$ 0.120 & 0.717 $\pm$ 0.120 & 0.711 $\pm$ 0.118 & 0.113 $\pm$ 0.111 & 0.112 $\pm$ 0.111 & 0.112 $\pm$ 0.128 & 0.110 $\pm$ 0.000 \\
reverse light off (masq.) & stealthy & 0.540 $\pm$ 0.183 & 0.537 $\pm$ 0.187 & 0.532 $\pm$ 0.178 & 0.000 $\pm$ 0.000 & 0.000 $\pm$ 0.000 & 0.001 $\pm$ 0.001 & 0.158 $\pm$ 0.000 \\
reverse light on (masq.) & stealthy & 0.096 $\pm$ 0.073 & 0.096 $\pm$ 0.072 & 0.089 $\pm$ 0.068 & 0.003 $\pm$ 0.004 & 0.002 $\pm$ 0.003 & 0.004 $\pm$ 0.004 & 0.018 $\pm$ 0.000 \\
max coolant temp$^\dagger$ & coarse & 0.047 $\pm$ 0.047 & 0.047 $\pm$ 0.047 & 0.047 $\pm$ 0.047 & 0.000 $\pm$ 0.000 & 0.000 $\pm$ 0.000 & 0.000 $\pm$ 0.000 & 0.093 $\pm$ 0.000 \\
max coolant temp (masq.)$^\dagger$ & stealthy & 0.085 $\pm$ 0.054 & 0.085 $\pm$ 0.054 & 0.085 $\pm$ 0.054 & 0.000 $\pm$ 0.000 & 0.000 $\pm$ 0.000 & 0.000 $\pm$ 0.000 & 0.140 $\pm$ 0.000 \\
pooled (primary pool) & -- & 0.545 $\pm$ 0.059 & 0.544 $\pm$ 0.060 & 0.538 $\pm$ 0.056 & 0.055 $\pm$ 0.041 & 0.054 $\pm$ 0.041 & 0.056 $\pm$ 0.047 & 0.276 $\pm$ 0.000 \\
\bottomrule
\end{tabular}
\end{table*}


\subsection{Per-attack-type quantization loss (RQ1)}\label{sec:res-rq1}

\textbf{Overall and per-type loss (H1).} H1 predicted that the pooled F1 loss from float32 to INT8 stays small but is
not spread evenly over the attack types. Table~\ref{tab:recall} gives the test recall per attack type, model and
quantization level (F1 in Table~\ref{tab:f1}; heat map of the recall loss in Fig.~\ref{fig:heatmap}). The largest
pooled INT8 F1 loss is \PooledDFOneMax{} (1D-CNN on ROAD), far below the bound of \HOneMax{} in the verdict rule. Of the
\NTypeRows{} rows of primary attack type and quantized model, \NLossRows{} have an INT8 recall-loss interval above
zero, \NZeroRows{} have an interval that contains zero, and \NGainRows{} lies below zero. The \NLossRows{} losses
all belong to the 1D-CNN on ROAD and amount to \LossRowsRange, less than one recall point. The single gain is the
conv-AE on Car-Hacking gear spoofing ($\dRecall{} =$ \GearGainCI, the same sign in every seed), whose $\Delta$ROC-AUC is
\GearGainAUC: the ranking of windows does not change, only the position of the threshold in the tail of the
calibration scores. H1 is therefore supported by its rule, but the unevenness it detects is negligible in size. On
Car-Hacking the 1D-CNN is at the recall ceiling (pooled recall \CHCNNRecall) and its $\dRecall$ is exactly zero for
every type and level, so this dataset leaves no room to observe a loss and contributes little to RQ1.

\textbf{Coarse versus stealthy attacks (H2).} H2 predicted that the INT8 loss is larger for stealthy than for coarse
attacks, and that coarse attacks (DoS, fuzzing) are still detected with high performance under INT8. The primary
evidence comes from fabrication--masquerade pairs of the same source capture, which carry the same injected payload
and differ only in the timing trace (Section~\ref{sec:threat}); \NPairUnits{} 1D-CNN pairs pass the floor rule
(Fig.~\ref{fig:delta}c, Table~\ref{tab:stats} block B). The INT8 loss on the two sides of a pair is of the same size
(\PairSideRange{} over the six sides), and the mean paired difference is
$D =$ \DRQTwoCI{} under INT8 and \DRQOneCI{} under float16. The interval of every single pair also contains zero. F1
leads to the same result ($D_{\mathrm{F1}} =$ \DFQTwoCI), and resampling windows independently instead of in blocks
barely changes the interval (\DRQTwoIid). With \NPairUnits{} pairs, a Wilcoxon signed-rank test cannot reach a
$p$-value below \PairMinP, so no test was run and the effect size and its interval are reported instead.

At the group level (Fig.~\ref{fig:delta}a--b, Table~\ref{tab:stats} block A), the 1D-CNN under INT8 has Cliff's
$\delta =$ \GroupCNNQTwoCI{} in the direction of H2, with exact Mann--Whitney $p = \GroupCNNQTwoP$ and Holm-adjusted
$p = \GroupCNNQTwoPHolm$. The effect size looks large, but the mean loss is \GroupMeanStealthy{} for the stealthy and
\GroupMeanCoarse{} for the coarse types, an absolute difference of \GroupDiffPts{} recall points, and the interval of
$\delta$ contains zero. Under float16 the 1D-CNN gives $\delta =$ \GroupCNNQOneCI{} ($p = \GroupCNNQOneP$, Holm
$p = \GroupCNNQOnePHolm$). For the conv-AE under INT8, $\delta =$ \GroupAEQTwoCI{} ($p = \GroupAEQTwoP$, Holm
$p = \GroupAEQTwoPHolm$); its recall on ROAD lies below the floor for most types (grey in Fig.~\ref{fig:delta}b), where
a loss is bounded by construction. On $\Delta$F1 the 1D-CNN under INT8 has $\delta =$ \GroupCNNQTwoFCI, and no
Holm-adjusted $p$ of either ROAD family is below \GroupMinHolm. With \GroupStealthy{} stealthy against \GroupCoarse{}
coarse types, the smallest attainable two-sided $p$ is \GroupMinPFrac{} $=$ \GroupMinP, which the Holm correction
raises to \GroupMinPHolm, above the conventional five-percent level: adjusted significance is unreachable in this
design whatever the data. Adding the Car-Hacking coarse types (\GroupStealthy{} against \SensCoarse) gives
$\delta =$ \SensCNNQTwoCI, $p = \SensCNNQTwoP$ and Holm $p = \SensCNNQTwoPHolm$, but the added types have a loss of
exactly zero because of the ceiling, so they carry the difference between the datasets rather than between the groups.

The second part of H2 also fails, for a reason unrelated to quantization (Table~\ref{tab:recall}): every coarse type keeps an INT8 recall of
at least \HTwoRecallMin{} (lowest value \CoarseOtherMin) except ROAD fuzzing for the conv-AE (\AEFuzzQTwo), whose
float32 recall is already \AEFuzzQZero{} $\pm$ \AEFuzzQZeroSD.

\emph{H2 is not supported.} INT8 produces a small recall loss for the 1D-CNN on ROAD, and this loss is of the same
size for a masquerade attack and for the fabrication attack it was derived from. In this design, the
stealthy-specific additional INT8 loss does not exceed about \DBoundPts{} recall points (upper end of the interval of
the mean $D$). A smaller stealthy-specific loss cannot be excluded, and the bound rests on \NPairUnits{} pairs of one
model on one dataset whose stealthy captures are synthetically derived (Section~\ref{sec:limitations}).

\begin{figure*}[t]
\centering
\includegraphics[width=\textwidth]{fig3_delta_recall_coarse_vs_stealthy}
\caption{Recall loss under quantization is small and not larger for stealthy attacks. Per-type
$\dRecall{} = \mathrm{Recall}(\Qfp) - \mathrm{Recall}(\mathrm{Q}k)$ with \CIPct{} hierarchical bootstrap intervals for
(a) the 1D-CNN and (b) the conv-AE on the ROAD primary types (Cliff's $\delta$ and Holm-adjusted $p$ of the group test
under each level), and (c) the paired masquerade-minus-fabrication difference $D$ for the 1D-CNN, whose INT8 mean is
\DRQTwoCI. Grey points have a float32 recall below \FloorRecall, where $\dRecall$ is uninformative.}
\label{fig:delta}
\end{figure*}

\begin{table*}[t]
\centering
\footnotesize
\setlength{\tabcolsep}{2.5pt}
\caption{Statistical comparisons ($\Delta$ = Q0 $-$ Qk; positive = loss) with 95\% hierarchical bootstrap intervals over seeds and test windows (moving blocks, B = 10,000). min. p: smallest p-value the design can attain. In block A it is 2/126 = 0.016, which the Holm correction raises to 0.063 > 0.05, so adjusted significance is unattainable whatever the data; block B is not tested (min. p 0.25). Block C: the pre-registered statistic, the mean over attack types of $|\Delta$AE$|$ $-$ $|\Delta$CNN$|$, is biased upward for the model whose $\Delta$ varies more across seeds; the corrected statistic is the signed difference $\Delta$AE $-$ $\Delta$CNN. The final H4 verdict (partially supported) rests on both, read together with the output-change measures.}
\label{tab:stats}
\begin{tabular}{@{}lllrrrr@{}}
\toprule
Model & Level & Statistic & Estimate [95\% CI] & $p$ & $p_\mathrm{Holm}$ & min. $p$ \\
\midrule
\multicolumn{7}{l}{\textit{A: coarse vs stealthy, group level, ROAD (4 stealthy vs 5 coarse types; exact Mann--Whitney U, Holm over 4 tests)}} \\
1D-CNN & Q1 & Cliff's $\delta$ of $\Delta$Recall & 0.65 [0.00, 0.75] & 0.135 & 0.444 & 0.0159 \\
1D-CNN & Q2 & Cliff's $\delta$ of $\Delta$Recall & 0.70 [$-$0.30, 0.90] & 0.111 & 0.444 & 0.0159 \\
conv-AE & Q1 & Cliff's $\delta$ of $\Delta$Recall & 0.20 [$-$0.20, 0.50] & 0.683 & 1 & 0.0159 \\
conv-AE & Q2 & Cliff's $\delta$ of $\Delta$Recall & $-$0.10 [$-$0.50, 0.50] & 0.905 & 1 & 0.0159 \\
1D-CNN & Q1 & Cliff's $\delta$ of $\Delta$F1 & 0.45 [0.00, 0.75] & 0.302 & 1 & 0.0159 \\
1D-CNN & Q2 & Cliff's $\delta$ of $\Delta$F1 & 0.10 [$-$0.30, 0.70] & 0.905 & 1 & 0.0159 \\
conv-AE & Q1 & Cliff's $\delta$ of $\Delta$F1 & 0.20 [$-$0.40, 0.50] & 0.683 & 1 & 0.0159 \\
conv-AE & Q2 & Cliff's $\delta$ of $\Delta$F1 & $-$0.10 [$-$0.50, 0.50] & 0.905 & 1 & 0.0159 \\
\midrule
\multicolumn{7}{l}{\textit{A': sensitivity, Car-Hacking coarse types added (4 vs 9; dataset and group confounded)}} \\
1D-CNN & Q1 & Cliff's $\delta$ of $\Delta$Recall & 0.69 [0.00, 0.75] & 0.0182 & 0.0615 & 0.0028 \\
1D-CNN & Q2 & Cliff's $\delta$ of $\Delta$Recall & 0.83 [$-$0.06, 0.94] & 0.0154 & 0.0615 & 0.0028 \\
conv-AE & Q1 & Cliff's $\delta$ of $\Delta$Recall & 0.17 [$-$0.22, 0.50] & 0.716 & 1 & 0.0028 \\
conv-AE & Q2 & Cliff's $\delta$ of $\Delta$Recall & $-$0.11 [$-$0.50, 0.44] & 0.796 & 1 & 0.0028 \\
1D-CNN & Q1 & Cliff's $\delta$ of $\Delta$F1 & 0.47 [$-$0.44, 0.69] & 0.199 & 0.794 & 0.0028 \\
1D-CNN & Q2 & Cliff's $\delta$ of $\Delta$F1 & 0.50 [0.19, 0.83] & 0.199 & 0.794 & 0.0028 \\
conv-AE & Q1 & Cliff's $\delta$ of $\Delta$F1 & 0.28 [$-$0.42, 0.50] & 0.449 & 0.898 & 0.0028 \\
conv-AE & Q2 & Cliff's $\delta$ of $\Delta$F1 & $-$0.11 [$-$0.56, 0.47] & 0.796 & 0.898 & 0.0028 \\
\midrule
\multicolumn{7}{l}{\textit{B: paired fabrication--masquerade units, 1D-CNN, ROAD (3 pairs; exploratory, not tested)}} \\
1D-CNN & Q1 & mean D ($\Delta$Recall) & 0.0010 [0.0000, 0.0029] & -- & -- & 0.25 \\
1D-CNN & Q2 & mean D ($\Delta$Recall) & 0.0006 [$-$0.0035, 0.0057] & -- & -- & 0.25 \\
1D-CNN & Q1 & mean D ($\Delta$F1) & 0.0010 [0.0000, 0.0030] & -- & -- & 0.25 \\
1D-CNN & Q2 & mean D ($\Delta$F1) & $-$0.0002 [$-$0.0033, 0.0035] & -- & -- & 0.25 \\
\midrule
\multicolumn{7}{l}{\textit{C: supervised vs unsupervised, conv-AE vs 1D-CNN (Wilcoxon signed-rank, Holm over 4 tests)}} \\
conv-AE vs 1D-CNN & Q1 & $|\Delta$AE$|$ $-$ $|\Delta$CNN$|$, ROC-AUC (pre-registered) & $8.4\times10^{-6}$ [$1.9\times10^{-6}$, $1.2\times10^{-4}$] & 0.893 & 1 & 0.000244 \\
conv-AE vs 1D-CNN & Q2 & $|\Delta$AE$|$ $-$ $|\Delta$CNN$|$, ROC-AUC (pre-registered) & 0.0062 [0.0048, 0.0272] & 0.000732 & 0.00293 & 0.000244 \\
conv-AE vs 1D-CNN & Q1 & $|\Delta$AE$|$ $-$ $|\Delta$CNN$|$, recall (pre-registered) & 0.0009 [$-$0.0003, 0.0048] & 0.5 & 1 & 0.0312 \\
conv-AE vs 1D-CNN & Q2 & $|\Delta$AE$|$ $-$ $|\Delta$CNN$|$, recall (pre-registered) & 0.0062 [0.0033, 0.0265] & 0.25 & 0.75 & 0.0312 \\
conv-AE vs 1D-CNN & Q1 & $\Delta$AE $-$ $\Delta$CNN, ROC-AUC (corrected) & $-7.8\times10^{-5}$ [$-1.6\times10^{-4}$, $-2.4\times10^{-7}$] & -- & -- & -- \\
conv-AE vs 1D-CNN & Q2 & $\Delta$AE $-$ $\Delta$CNN, ROC-AUC (corrected) & 0.0025 [$-$0.0041, 0.0088] & -- & -- & -- \\
conv-AE vs 1D-CNN & Q1 & $\Delta$AE $-$ $\Delta$CNN, recall (corrected) & $-$0.0008 [$-$0.0048, 0.0004] & -- & -- & -- \\
conv-AE vs 1D-CNN & Q2 & $\Delta$AE $-$ $\Delta$CNN, recall (corrected) & $-$0.0002 [$-$0.0176, 0.0172] & -- & -- & -- \\
\bottomrule
\end{tabular}
\end{table*}


\subsection{Latency, size and recall (RQ2)}\label{sec:res-rq2}

H3 predicted that float16 is the most balanced point between loss and gain, and that INT8 gives the largest gain but
a marked loss on stealthy attacks. Table~\ref{tab:resources} lists the resources and Fig.~\ref{fig:pareto} the
trade-off.

\textbf{Size.} The float16 file is \SizeRatioQOne{} and the INT8 file \SizeRatioQTwo{} of the float32 file for both
deep models (1D-CNN \SizeCNNQZero{}, \SizeCNNQOne{} and \SizeCNNQTwo~MB; conv-AE \SizeAEQZero{}, \SizeAEQOne{} and
\SizeAEQTwo~MB). The XGBoost model takes \SizeXGBCH~MB (Car-Hacking) and \SizeXGBROAD~MB (ROAD).

\textbf{Latency.} Absolute latencies on the host are read as orders of magnitude only: \DeepLatRange~ms per window for
the deep models and \XGBLatRange~ms for XGBoost. Because the two measurement sessions disagree beyond the tolerance,
the direction of an effect is taken from within-session ratios and stated only where both sessions agree. INT8
lowers the median latency of the 1D-CNN in both sessions and on both datasets (ratio to float32 \CNNRatio, i.e.
\CNNCutPct\% lower), but not that of the conv-AE (\AERatio). The conv-AE is \AEvsCNN{} times and XGBoost \XGBvsCNN{}
times slower than the float32 1D-CNN. The float16 models have no CPU latency, because the CPU provider has no float16
kernels for their layers (Section~\ref{sec:method}); the tail percentiles in Table~\ref{tab:resources} are not
interpreted.

\textbf{Pareto front and budgets (Fig.~\ref{fig:pareto}).} On ROAD, with the recall over the stealthy types as accuracy axis, the robust
front holds only the float32 1D-CNN (\StealthCNNQZero) and the INT8 1D-CNN (\StealthCNNQTwo): INT8 buys a
\CNNCutPct\% lower latency and a file of \SizeRatioQTwo{} times the size for a stealthy recall that is \StealthDiff{}
lower, with largely overlapping intervals. The conv-AE (\StealthAEQZero{} and \StealthAEQTwo) and XGBoost
(\StealthXGB) are dominated. On Car-Hacking, with the pooled recall, the front holds the INT8 1D-CNN
(\PooledCHCNNQTwo) and XGBoost (\PooledCHXGB), which is an order of magnitude slower and several times larger. On the
size--recall front, where float16 can enter, all three 1D-CNN levels lie on the ROAD front, float16 between the other
two (\SizeCNNQOne~MB, stealthy recall \StealthCNNQOne). Every candidate stays far below the window budget
(\BudgetWindowCH~ms on Car-Hacking, \BudgetWindowROAD~ms on ROAD). Under the strict budget (\BudgetStrictCH~ms and
\BudgetStrictROAD~ms) every deep model fits (at most \DeepLatMax~ms), whereas XGBoost misses it on ROAD
(\XGBSlowROAD~ms) and fits on Car-Hacking (\XGBSlowCH~ms) by a margin below the precision that the host measurement
supports. Neither budget therefore separates the deep candidates on this host, and on ROAD the INT8 1D-CNN is the most
suitable combination: the smallest and fastest candidate, with a stealthy recall within the interval of its float32
version. These are host measurements; compliance with a budget here is not evidence for an ECU.

\emph{H3 is partially supported.} The float16 models lose less recall than the INT8 ones (1D-CNN on ROAD: mean $|\dRecall|$
\AbsDRQOne{} against \AbsDRQTwo) and halves the file, but its latency gain could not be measured, and INT8 does not
cause the predicted stealthy-specific loss (H2).

\begin{figure*}[t]
\centering
\includegraphics[width=\textwidth]{fig5_pareto}
\caption{Latency--recall trade-off on the host CPU: (a) recall over the ROAD stealthy types (\CIPct{} interval),
(b) Car-Hacking pooled recall ($\pm$ one s.d. over seeds), against the median per-window latency. The horizontal bar is
the range of the two measurement sessions and the marker sits at the slower session; marker area is file size; the
dotted line is the robust Pareto front; dashed and dash-dotted lines are the strict (one-message) and window latency
budgets.}
\label{fig:pareto}
\end{figure*}

\begin{table*}[t]
\centering
\footnotesize
\setlength{\tabcolsep}{3pt}
\caption{Resources on the host CPU (4 threads, batch 1, seed-0 models). Latency is the range over two measurement sessions; p95 is reported but not interpreted, absolute latencies are orders of magnitude only. Q1 has no CPU latency (no fp16 kernel in the CPU provider). A model meets a budget if the slower session's median is below it.}
\label{tab:resources}
\begin{tabular}{@{}llrrrrrrcc@{}}
\toprule
& & \multicolumn{2}{c}{File} & \multicolumn{2}{c}{Latency (ms), 2 sessions} & Ratio vs & Peak RSS & \multicolumn{2}{c}{Meets budget} \\
\cmidrule(lr){3-4}\cmidrule(lr){5-6}\cmidrule(lr){9-10}
Model & Level & MB & /Q0 & median & p95 & CNN Q0 & (MB) & 1 msg & window \\
\midrule
\multicolumn{10}{l}{\textit{Car-Hacking (budgets: 1 message 0.526 ms, window 52.6 ms)}} \\
1D-CNN & Q0 (fp32) & 0.103 & 1.00 & 0.040--0.042 & 0.054--0.055 & 1.00--1.00 & 65.2--65.4 & yes & yes \\
1D-CNN & Q1 (fp16) & 0.054 & 0.52 & -- & -- & -- & -- & -- & -- \\
1D-CNN & Q2 (INT8) & 0.039 & 0.38 & 0.029--0.030 & 0.034--0.039 & 0.70--0.74 & 65.6--65.7 & yes & yes \\
conv-AE & Q0 (fp32) & 0.111 & 1.00 & 0.102--0.110 & 0.110--0.192 & 2.42--2.73 & 64.6--64.7 & yes & yes \\
conv-AE & Q1 (fp16) & 0.058 & 0.52 & -- & -- & -- & -- & -- & -- \\
conv-AE & Q2 (INT8) & 0.042 & 0.38 & 0.110--0.116 & 0.161--0.163 & 2.61--2.88 & 65.3--65.4 & yes & yes \\
XGBoost & native & 0.205 & -- & 0.489--0.518 & 0.952--1.274 & 12.13--12.30 & 269.1--269.2 & yes & yes \\
\midrule
\multicolumn{10}{l}{\textit{ROAD (budgets: 1 message 0.419 ms, window 41.9 ms)}} \\
1D-CNN & Q0 (fp32) & 0.103 & 1.00 & 0.040--0.046 & 0.045--0.058 & 1.00--1.00 & 65.0--65.2 & yes & yes \\
1D-CNN & Q1 (fp16) & 0.054 & 0.52 & -- & -- & -- & -- & -- & -- \\
1D-CNN & Q2 (INT8) & 0.039 & 0.38 & 0.030--0.030 & 0.033--0.034 & 0.64--0.75 & 65.5--65.8 & yes & yes \\
conv-AE & Q0 (fp32) & 0.111 & 1.00 & 0.102--0.107 & 0.110--0.126 & 2.30--2.57 & 64.7--64.8 & yes & yes \\
conv-AE & Q1 (fp16) & 0.058 & 0.52 & -- & -- & -- & -- & -- & -- \\
conv-AE & Q2 (INT8) & 0.042 & 0.38 & 0.107--0.108 & 0.113--0.172 & 2.32--2.69 & 64.9--65.0 & yes & yes \\
XGBoost & native & 0.274 & -- & 0.484--0.500 & 0.937--1.238 & 10.40--12.57 & 269.6--269.8 & no & yes \\
\bottomrule
\end{tabular}
\end{table*}


\subsection{Supervised versus unsupervised models (RQ3)}\label{sec:res-h4}

H4 predicted that the reconstruction-based unsupervised model is more sensitive to quantization than the supervised
classifier, because its anomaly score relies on small differences of reconstruction error. The verdict went through
three stages, which are reported together (Table~\ref{tab:stats} block C; Fig.~\ref{fig:h4}).

\textbf{First stage: pre-registered statistic.} The rule compares, per attack type, the absolute losses of the two models,
$d_t = |\Delta_{\mathrm{AE},t}| - |\Delta_{\mathrm{CNN},t}|$, averaged over types. Under INT8 the mean $d_t$ is
\PreAUCQTwoCI{} on $\Delta$ROC-AUC over \PreAUCUnits{} types (Wilcoxon $p = \PreAUCQTwoP$, Holm $p = \PreAUCQTwoPHolm$;
Cliff's $\delta$ of the absolute losses \PreAUCQTwoDelta) and \PreRecQTwoCI{} on $\dRecall$ over the \PreRecUnits{}
types where both models reach a float32 recall of at least \RecallFloorAS{} ($p = \PreRecQTwoP$, Holm
$p = \PreRecQTwoPHolm$; $\delta = \PreRecQTwoDelta$; with \PreRecUnits{} types the smallest attainable $p$ is
\PreRecMinP, \PreRecMinPHolm{} after Holm). Under float16 the statistic is \PreAUCQOneCI{} on $\Delta$ROC-AUC
($p = \PreAUCQOneP$). Since both INT8 intervals lie above zero, the rule declares H4 supported.

\textbf{Second stage: corrected statistic.} This verdict cannot stand alone. The expected absolute value of a noisy estimate
is positive even when the true effect is zero, and it grows with the noise. The losses of the conv-AE vary far more
across seeds than those of the 1D-CNN (seed s.d. of the INT8 $\Delta$ROC-AUC on ROAD correlated signal: \SDAECorr{}
against \SDCNNCorr), so $d_t$ is biased upward in favour of H4; the point estimate near the lower end of its interval
is the symptom, and the Wilcoxon test on the same $d_t$ inherits the bias. The corrected statistic is the signed
difference $\Delta_{\mathrm{AE}} - \Delta_{\mathrm{CNN}}$, which has no such bias: \SgnAUCQTwoCI{} on $\Delta$ROC-AUC
and \SgnRecQTwoCI{} on $\dRecall$ under INT8. Neither model's own mean $\Delta$ROC-AUC differs from zero (conv-AE
\MeanAEAUC, 1D-CNN \MeanCNNAUC). Under float16 the signed difference, \SgnAUCQOneCI, excludes zero but is negligible
and favours the conv-AE. The directional reading of H4, that the unsupervised model loses more performance, is
therefore not supported. Measures of output change, which are direct measurements rather than differences of two
noisy estimates, do separate the models (Fig.~\ref{fig:h4}c--d): on ROAD, INT8 changes the decision for \ChurnAEROAD{} (seeds
\ChurnAEROADRange) of the attack windows of the conv-AE against \ChurnCNNROAD{} (\ChurnCNNROADRange) for the 1D-CNN,
\ChurnRatioROAD{} times as many, and the Spearman correlation between float32 and INT8 window scores is \SpearAEROAD{}
(\SpearAEROADRange) against \SpearCNNROAD. The changed decisions go both ways, so they do not add up to a net loss. On
Car-Hacking the two models are similar (Spearman \SpearAECH{} and \SpearCNNCH).

\textbf{Third stage: final verdict.} \emph{H4 is partially supported.} Supported: INT8 changes the output of the
unsupervised model more than that of the supervised one (score ranks and decisions on ROAD, and the calibration shift
of Section~\ref{sec:res-calib}). Not supported: once the threshold is re-fitted on the quantized model's own
validation scores, the unsupervised model does not lose more recall or ROC-AUC. On timing: the verdict rule was written
after the descriptive results and before the intervals; the corrected statistic and the output-change measures were
added after a smoke run of the analysis, before the full run; and the final verdict was decided by the authors after
all results had been seen. The verdict of the pre-registered rule is reported, not replaced.

\subsection{Calibration shift}\label{sec:res-calib}

If the float32 threshold is kept for the quantized model instead of being re-fitted (Fig.~\ref{fig:calib}), INT8
shifts the reconstruction error of the conv-AE on ROAD: the validation false-positive rate rises from \CalValQZero{} to
\CalValQTwo{} (seeds \CalValRange), a \CalFold-fold increase depending on the seed, and the false-positive rate on the
attack-free test captures from \CalTestQZero{} to \CalTestQTwo{} (\CalTestRange). The pooled recall rises from
\CalRecQZero{} to \CalRecQTwo{} (\CalRecRange), but this is not better detection: it comes from the false alarms, and
in this mode the per-type recall change reaches \CalTypeMin. On Car-Hacking the shift of the conv-AE is small
(validation false-positive rate \CalCHValQZero{} to \CalCHValQTwo), and the 1D-CNN hardly shifts at all (largest
change of a mean false-positive rate \ShiftCNNMax); float16 does not shift either model (at most \ShiftQOneMax). A
quantized autoencoder must therefore be re-thresholded on its own validation scores. The primary analysis does this,
which is why the shift does not appear in RQ1 and RQ3; it is the deployment-side face of the output sensitivity found
for H4.

\begin{figure*}[t]
\centering
\includegraphics[width=\textwidth]{fig8_calibration_shift}
\caption{Calibration shift when the float32 threshold is kept for the quantized model: attack-free (a) validation and
(b) test false-positive rate and (c) pooled test recall, one line per seed. INT8 moves the reconstruction error of the
ROAD conv-AE so far that the validation false-positive rate rises \CalFold-fold and the apparent recall gain is false
alarms, while float16 and the 1D-CNN do not shift.}
\label{fig:calib}
\end{figure*}

\subsection{Further observations}\label{sec:res-other}

Masquerade attacks are not harder than fabrication for the 1D-CNN: its float32 recall is higher on the masquerade
side in two of the three pairs (max speedometer \RecCNNSpeedMasq{} against \RecCNNSpeedFab, reverse light off
\RecCNNRevMasq{} against \RecCNNRevFab; correlated signal \RecCNNCorrMasq{} against \RecCNNCorrFab). For XGBoost,
which sees window summary features, fabrication is clearly easier (max speedometer \RecXGBSpeedFab{} against
\RecXGBSpeedMasq, reverse light off \RecXGBRevFab{} against \RecXGBRevMasq; Table~\ref{tab:recall}). A plausible but
untested reading is that XGBoost relies on the timing trace that masquerade removes, whereas the 1D-CNN reads the
payload bytes. INT8 brings no latency gain for the conv-AE (\AERatio; Table~\ref{tab:resources}), which is consistent with its \HeavySimAE{}
transposed-convolution layers, out of \HeavyAE{} heavy layers, running as simulated quantization without an INT8
kernel (Section~\ref{sec:method}); this explanation was not tested separately.
